%%%%%%%%%%%%%%%%%
% This is an sample CV template created using altacv.cls
% (v1.6, 21 May 2021) written by LianTze Lim (liantze@gmail.com). Now compiles with pdfLaTeX, XeLaTeX and LuaLaTeX.
%
%% It may be distributed and/or modified under the
%% conditions of the LaTeX Project Public License, either version 1.3
%% of this license or (at your option) any later version.
%% The latest version of this license is in
%%    http://www.latex-project.org/lppl.txt
%% and version 1.3 or later is part of all distributions of LaTeX
%% version 2003/12/01 or later.
%%%%%%%%%%%%%%%%

%% Use the "normalphoto" option if you want a normal photo instead of cropped to a circle
% \documentclass[10pt,a4paper,normalphoto]{altacv}

\documentclass[10pt,a4paper,ragged2e,withhyper]{altacv}
%% AltaCV uses the fontawesome5 and packages.
%% See http://texdoc.net/pkg/fontawesome5 for full list of symbols.

% Change the page layout if you need to
\geometry{left=1.25cm,right=1.25cm,top=1.4cm,bottom=1.4cm,columnsep=0.8cm}

% The paracol package lets you typeset columns of text in parallel
\usepackage{paracol}

% Change the font if you want to, depending on whether
% you're using pdflatex or xelatex/lualatex

% If using pdflatex:
% \usepackage[rm]{roboto}
% \usepackage[defaultsans]{lato}
% \usepackage{tgadventor}
% \usepackage[sfdefault]{AlegreyaSans}
\usepackage{cmbright}
% \usepackage[sfdefault]{biolinum}
\usepackage{pifont}
% \usepackage{sourcesanspro}
\renewcommand{\familydefault}{\sfdefault}
\usepackage{academicons}
\usepackage{relsize}
% Change the colours if you want to
\definecolor{SlateGrey}{HTML}{2E2E2E}
\definecolor{LightGrey}{HTML}{666666}
\definecolor{LightestGrey}{HTML}{D3D3D3}
\definecolor{DarkPastelRed}{HTML}{450808}
\definecolor{PastelRed}{HTML}{8F0D0D}
\definecolor{GoldenEarth}{HTML}{E7D192}
\definecolor{CoolBlack}{HTML}{002E63}
\definecolor{OxfordBlue}{HTML}{4771AD}
\colorlet{name}{black}
\colorlet{personalinfo}{OxfordBlue}
\colorlet{heading}{OxfordBlue}
\colorlet{headingrule}{LightestGrey}
\colorlet{accent}{black}
\colorlet{emphasis}{SlateGrey}
\colorlet{body}{black}
%\usepackage[colorlinks,linkcolor=OxfordBlue,citecolor=red,urlcolor=blue,bookmarks=false,hypertexnames=true]{hyperref}

\makeatletter
\newcommand\footnoteref[1]{\protected@xdef\@thefnmark{\ref{#1}}\@footnotemark}
\makeatother

% Change some fonts, if necessary
% \renewcommand{\namefont}{\LARGE\bfseries}
% \renewcommand{\personalinfofont}{\footnotesize}
% \renewcommand{\cvsectionfont}{\LARGE\rmfamily\bfseries}
% \renewcommand{\cvsubsectionfont}{\large\bfseries}

% Change the bullets for itemize and rating marker
% for \cvskill if you want to
\renewcommand{\itemmarker}{{\small\textbullet}}
\renewcommand{\ratingmarker}{\faCircle}
%% Use (and optionally edit if necessary) this .tex if you
%% want to use an author-year reference style like APA(6)
%% for your publication list
% When using APA6 if you need more author names to be listed
% because you're e.g. the 12th author, add apamaxprtauth=12
\usepackage[backend=biber,style=apa6,sorting=ydnt]{biblatex}
\defbibheading{pubtype}{\cvsubsection{#1}}
\renewcommand{\bibsetup}{\vspace*{-\baselineskip}}
\AtEveryBibitem{\makebox[\bibhang][l]{\itemmarker}}
\setlength{\bibitemsep}{0.25\baselineskip}
\setlength{\bibhang}{1.25em}


%% Use (and optionally edit if necessary) this .tex if you
%% want an originally numerical reference style like IEEE
%% for your publication list
% \usepackage[backend=biber,style=ieee,sorting=ydnt]{biblatex}
%% For removing numbering entirely when using a numeric style
\setlength{\bibhang}{1.25em}
\DeclareFieldFormat{labelnumberwidth}{\makebox[\bibhang][l]{\itemmarker}}
\setlength{\biblabelsep}{0pt}
\defbibheading{pubtype}{\cvsubsection{#1}}
\renewcommand{\bibsetup}{\vspace*{-\baselineskip}}


%% sample.bib contains your publications
\addbibresource{sample.bib}
\begin{document}
\name{Neelkamal Somisetty}
% \tagline{Your Position or Tagline Here}

\personalinfo{
  \tag{\faEnvelope}\info{\href{mailto:neelkamalsomisetty96@gmail.com} {neelkamalsomisetty96@gmail.com}}
%   \homepage{}
%   \linkedin{LinkedIn}{linkedin.com/in/vishal-mohanty}
  \printinfo{\faLinkedin}{LinkedIn}[https://www.linkedin.com/in/neelkamal-somisetty-ph-d-1b973911b/]
  \printinfo{\faGithubSquare}{GitHub}[https://github.com/Neelkamal18]
  \printinfo{\faGoogle}{Google Scholar}[https://scholar.google.com/citations?hl=en]
  \phone{+1-979-777-2845}
%   \github{your_id}
%   \orcid{0000-0000-0000-0000}
}

\makecvheader
%% Depending on your tastes, you may want to make fonts of itemize environments slightly smaller
% \AtBeginEnvironment{itemize}{\small}

%% Set the left/right column width ratio to 6:4.
\columnratio{0.341}

% Start a 2-column paracol. Both the left and right columns will automatically
% break across pages if things get too long.
\begin{paracol}{2}
\vspace{-0.5em}
\cvsection{\large Education}
\vspace{-0.25em}

\cvproject{\small Texas A\&M University}{\small Ph.D. Mechanical Engineering \hspace{0.07em} GPA: 4.0/4.0}{\small 2021-2024}{\small Advisor: {\color{OxfordBlue}\href{https://engineering.tamu.edu/mechanical/profiles/darbha-swaroop.html}{Dr.Swaroop Darbha}}}

\cvproject{\small Texas A\&M University}{\small M.S. Mechanical Engineering}{\small 2018-2021}{}

\cvproject{\small IIT Madras}{\small B.Tech Mechanical Engineering}{\small 2014-2018}{}

\vspace{-0.5em}
\divider
\vspace{-0.5em}

\cvsection{\large Skills}

\vspace{-0.75em}
\cvevent{\small Languages \& Software Tools}{}{\vspace{-0.25em}\color{black} Python, Julia, C++, OpenCV, ROS, MATLAB, Gurobi, CPLEX, Git, Gazebo, PyBullet, Docker, Streamlit, Linux}
\vspace{-0.5em}
\cvevent{\small ML \& Data}{}{\vspace{-0.25em}\color{black} PyTorch, TensorFlow, Keras, CUDA, Scikit-learn, BigQuery, SQL, Dataflow, Stable-Baselines3, OpenAI Gymnasium}

\vspace{-0.5em}
\divider
\vspace{-0.5em}

\cvsection{\large Selected Courses}
\vspace{-0.25em}

\color{black} \textbf{Modern Computer vision}, \textbf{Robotic software and perception}, \textbf{Machine learning}, \textbf{Sensor fusion}, Data structures and algorithms, Combinatorial and Convex optimization, Spectral Graph Theory, Control systems, Statistics

\vspace{-0.25em}
\divider
\vspace{-0.5em}

\cvsection{\large Publications}
\vspace{-0.25em}


\begin{itemize}
    \item \textbf{Somisetty, N.}, et al., (2024, Jul). {\color{OxfordBlue}\href{https://ieeexplore.ieee.org/abstract/document/10605101}{{Optimal Robust Network Design: Formulations and Algorithms for Maximizing Algebraic Connectivity}}}. IEEE TCNS
    \vspace{-0.3em}
    \item \textbf{Somisetty, N.}, et al., (2024, Dec). {\color{OxfordBlue}\href{https://arxiv.org/abs/2409.15506}{Spectral Graph Theoretic Methods for Enhancing Network Robustness in Robot Localization}}. IEEE CDC
    \vspace{-0.3em}
    \item \textbf{Somisetty, N.}, et al., (2025). Optimal Spectral Sparsification for Pose-Graph SLAM by Maximizing Algebraic Connectivity. IEEE RAL
    \vspace{-0.3em}
    \item \textbf{Somisetty, N.}, et al., (2025). Heuristic Sparsification for Pose-Graph SLAM by Maximizing Algebraic Connectivity. IEEE ICRA
    \vspace{-0.3em}
    \item \textbf{Somisetty, N.} and Darbha, S. (2024, Dec). {\color{OxfordBlue}\href{https://arxiv.org/abs/2411.07540}{Lateral String Stability in Autonomous and Connected Vehicle Platoons}}. IEEE ICVES
    \vspace{-0.3em}
    \item \textbf{Somisetty, N.}, et al., (2023, Sept). {\color{OxfordBlue}\href{https://ieeexplore.ieee.org/abstract/document/10422222}{Impact of Ordering of Heterogeneous Vehicles on String Stability.}} IEEE ITSC
    \vspace{-0.3em}
    \item  Liu, M., \textbf{Somisetty, N.}, et al., (2025). {\color{OxfordBlue}\href{https://arxiv.org/abs/2011.03587}{Robust Lateral Control of a Convoy of Autonomous and Connected Vehicles with Limited Preview}}. IEEE ITS
    \vspace{-0.25em}
\end{itemize}

% \divider

%% Switch to the right column. This will now automatically move to the second
%% page if the content is too long.
\switchcolumn
\vspace{-0.45em}
\cvsection{\large Professional Experience}
\vspace{-0.5em}

\cvproject{\small Autonomy Intern}{\href{https://www.monarchtractor.com/}{\small Monarch Tractor}}{\small September 2023 - May 2024}{ \small
\begin{itemize}
    \item \textbf{Trained} and evaluated the \textbf{SuperPoint CNN} using PyTorch on real and synthetic farm datasets, achieving efficient feature matching compared to SIFT and ORB
    \vspace{-0.4em}
    \item Designed a GPS-independent localization library using ROS for \textbf{sensor fusion of IMU, LiDAR, and visual odometry,} achieving precise positioning
    \vspace{-0.4em}
    \item Developed a bundle adjustment pipeline with a baseline constraint for a multi-view stereo, \textbf{enhancing 3D reconstruction accuracy by 11\% }in SfM
\end{itemize}}
\vspace{-0.5em}

\cvproject{\small Research Intern}{\small Applied Mathematics and Plasma Physics, Los Alamos National Laboratory}{\small May 2023 - August 2023}{\color{red}
\begin{itemize}
    \item Trained and evaluated the \textbf{SuperPoint neural network} using PyTorch on farm datasets, achieving efficient feature matching compared to SIFT
    \vspace{-0.4em}
    \item Designed a GPS-independent localization library using ROS for sensor fusion of IMU, LiDAR, and visual odometry, achieving precise positioning  
 \end{itemize}}
\vspace{-0.5em}

\cvproject{\small Graduate Research Assistant}{{\href{https://autonomy-its.engr.tamu.edu}{\small Systems, Controls and Applied Optimization Laboratory}}}{\small January  2020 - May 2023}{
\begin{itemize}
    \item Integrated research in \textbf{robotics, optimization, AI, and spectral graph theory} 
    \vspace{-0.4em}
    \item Designed formulations and algorithms to represent mathematical programs and optimize complex physical systems, including \textbf{robust network design} challenges
    \vspace{-0.4em}
    \item Utilized convex and discrete optimization techniques to \textbf{sparsify large-scale networks} for \textbf{UAV cooperative localization} and \textbf{pose-graph SLAM} applications
    \vspace{-0.4em}
    \item Developed an \textbf{optimality-based} \textbf{cutting-plane algorithm} for a \textbf{mixed-integer semidefinite program}, efficiently solving small-scale NP-complete problems
    \vspace{-0.4em}
    \item Devised \textbf{heuristic methods} for \textbf{enhancing network resilience} in UAV cooperative localization, and contributed to the development of \href{https://github.com/harshangrjn/LaplacianOpt.jl}{\color{OxfordBlue}\textbf{LaplacianOpt.jl}}, a Julia-based package for network analysis
    \vspace{-0.4em}
    \item Developed a \textbf{neighborhood search heuristic} for SLAM pose graph sparsification, cutting memory usage by \textbf{15\%} and boosting computational efficiency by \textbf{23\%}
\end{itemize}}


\vspace{-0.75em}
\divider
\vspace{-0.5em}

\cvsection{\large Selected Projects}
\vspace{-0.25em}

\cvproject{\small Natural Language-Guided Robotic Assistant}{}{\small 2024}{
    \begin{itemize}
        \item Designed a multi-modal robotic system integrating \textbf{VLMs}, \textbf{RL}, and the \textbf{BERT NLP model} to interpret and execute language-based instructions
        \vspace{-0.4em}
        \item Developed a \textbf{custom OpenAI Gym}-compatible environment to train \textbf{RL policies (PPO)} leveraging VLMs for \textbf{vision-language grounding} and \textbf{task planning}
        \vspace{-0.4em}
        \item Implemented a robust framework combining \textbf{CLIP} with ROS, enabling robots to perform dynamic tasks like \textbf{object manipulation and navigation}
    \end{itemize}
}

\vspace{-0.5em}
\cvproject{\small Cycle-GAN for Sim-to-Real Domain Adaptation}{}{\small 2024}{
    \begin{itemize}
        \item Developed a \textbf{CycleGAN} model to translate simulated robotic environments into realistic images, enhancing the \textbf{generalization of robotic perception systems}
        \vspace{-0.4em}
        \item Improved task performance by bridging the \textbf{sim-to-real gap}, enabling robotic systems trained in simulation to perform effectively in real-world environments
    \end{itemize}
}

\vspace{-0.5em}
\cvproject{\small Pose-Graph Sparsification for Multi-Session SLAM in UGVs}{}{\hspace{-0.5em}\small 2023}{
    \begin{itemize}
        \item Processed LiDAR data from multiple sessions using \textbf{Fast-LIO2 for odometry}, and utilized \textbf{ScanContext} and \textbf{DiSCo-SLAM} for detecting \textbf{loop closures}
        \vspace{-0.4em}
        \item Performed offline processing to construct and sparsify pose graphs using heuristics developed in \href{https://github.com/harshangrjn/LaplacianOpt.jl}{\color{OxfordBlue}\textbf{LaplacianOpt.jl}}
        \vspace{-0.4em}
        \item Conducted qualitative analysis with \textbf{LiDAR point clouds} and \textbf{GPS visualizations}, and quantitatively evaluated sparsifications using \textbf{GTSAM} and \textbf{SE-Sync} libraries 
    \end{itemize}
}
\vspace{-0.5em}
\cvproject{\small Deep Reinforcement Learning for Supply Chain}{}{\small 2023}{
    \begin{itemize}
        \item Reformulated a mixed-integer linear program as a Markov Decision Process and built a \textbf{custom OpenAI Gymnasium environment} to \textbf{test on-policy} and \textbf{off-policy algorithms}
        \vspace{-0.4em}
        \item Evaluated DRL approaches, including \textbf{PPO} and \textbf{DDPG}, against integer programming, demonstrating its
\textbf{effectiveness in handling dynamic constraints}
    \end{itemize}
}
\vspace{-0.5em}
\cvproject{\small Path Planning for KUKA Robotic Arm}{}{\small 2023}{
    \begin{itemize}
        \item Designed a \textbf{grid-based $A^*$ search algorithm} for path planning through maze and simulated the \textbf{robotic arm’s motion in ROS and Gazebo} for real-time testing
        \vspace{-0.4em}
        \item Optimized algorithms by integrating heuristic enhancements, \textbf{bidirectional search strategies,} and \textbf{memory-efficient data structures}, achieving real-time planning
    \end{itemize}
}
\vspace{-0.5em}
\cvproject{\small Sentiment Analysis for YELP reviews using BERT}{}{\small 2022}{
    \begin{itemize}
        \item Built a {\color{OxfordBlue}\href{https://github.com/Neelkamal18/Sentiment_analysis_for_YELP_reviews}{sentiment analysis pipeline}} using \textbf{transformer-based models}, incorporating techniques like \textbf{tokenization} and \textbf{attention mechanisms} to classify customer feedback with high accuracy and contextual understanding
    \end{itemize}
}
\vspace{-0.5em}
\cvproject{\small RAG for Knowledge-Enhanced NLP Applications}{}{\small 2024}{
    \begin{itemize}
        \item Leveraged \textbf{RAG} to integrate external knowledge into \textbf{NLP workflows}, enhancing accuracy and contextual relevance in query resolution
        \vspace{-0.4em}
        \item Built a scalable \textbf{RAG pipeline} combining \textbf{pre-trained language models} with real-time \textbf{retrieval systems} to deliver efficient and knowledge-driven text generation
    \end{itemize}
}



%\cvproject{Lateral Feedback Controller for Vehicles}{Advisor: Dr.Swaroop Darbha}{2019}{
%    \begin{itemize}
%        \item 
%    \end{itemize}
%}

%\cvproject{\href{https://www.raftarformularacing.com}{\small RAFTAR: FORMULA STUDENT TEAM}}{}{\small 2015-2017}
%{
%    \begin{itemize}
%        \item 
%    \end{itemize}
%}

%\divider

%\cvsection{Notable Awards}

%\begin{itemize}
%    \item Department of Energy graduate research fellowship (2020 - 2024).
%    \item Professor Rayudu Memorial fellowship at IIT Madras (2014 - 2018).
%    \item KVPY Scholar (2014) - awarded by the Government of India.
%\end{itemize}

%\cvsection{Publications}

%\nocite{*}

%\printbibliography[heading=none,title={\printinfo{\faUsers}{Conference Proceedings}},type=inproceedings]

%\printbibliography[heading=none,title={\printinfo{\faFile*[regular]}{Journal Articles}},type=article]

% \begin{itemize}
%     \item Auditing Indian Elections, \href{https://www.springer.com/gp/book/9783030306243}{\color{Red}LNCS Electronic Voting 2019} 
%     \item Threshold implementation of AES S-Box on FPGA,  \href{https://ieeexplore.ieee.org/abstract/document/8641387}{\color{Red}IEEE Embedded Systems Letters 2019}
% \end{itemize}

\end{paracol}


\end{document}